\begin{center}
\hypertarget{native-063-article046}{}\pdfbookmark[2]{46.}{article-046}46.
\end{center}
Motum hyperbolicum quoque aliquot exemplis illustravisse haud inutile erit, ad quem finem numeros artt. 23, 26 resumemus.

I. Data sunt $e=1{,}261\,8820$, $\log q=0{,}020\,1657$, $v=18^0\,51^{\prime}\,0^{\prime\prime}$: quaeritur $t$. Habemus
\begin{center}
\begin{tabular}{@{}ll@{}}
\makebox[20em][l]{$2\log\operatorname{tang}\tfrac{1}{2}v$\hfill .\,.\,.\,.\,.\,.\,.\,.} & $8{,}440\,2018$\\[.3em]
\makebox[20em][l]{$\log\dfrac{e-1}{e+1}$\hfill .\,.\,.\,.\,.\,.\,.\,.\,.\,.\,.\,.} & $9{,}063\,6357$\\[.4em]
\hline
\makebox[20em][l]{$\log T$\hfill .\,.\,.\,.\,.\,.\,.\,.\,.\,.\,.\,.\,.\,.} & $7{,}503\,8375$\\[.3em]
$T$\hfill $=$ & $0{,}003\,1903\,4$\\[.3em]
$\log B$\hfill $=$ & $0{,}000\,0001$\\[.3em]
$C$\hfill $=$ & $0{,}000\,0005$
\end{tabular}
\end{center}
\begin{center}
\begin{tabular}{@{}lrl l@{}}
\multicolumn{3}{@{}l}{\makebox[22em][l]{$\log\dfrac{2Bq^{\frac{3}{2}}}{k\BareRoot{(e-1)}}$\hfill .\,.\,.\,.\,.\,.\,.\,.\,.\,.}} & $2{,}386\,6444$\\[.4em]
\multicolumn{3}{@{}l}{\makebox[22em][l]{$\log A^{\frac{1}{2}}$\hfill .\,.\,.\,.\,.\,.\,.\,.\,.\,.\,.\,.\,.}} & $8{,}752\,4738$\\[.3em]
\hline
$\log$ & $13{,}77584$ & $=$ & $1{,}139\,1182$\\[.3em]
 & $0{,}13861$ & & \\
\cline{2-2}
 & $13{,}91445$ & $=t$ &
\end{tabular}
\end{center}
\begin{center}
\begin{tabular}{@{}ll@{}}
\makebox[22em][l]{$\log T$\hfill .\,.\,.\,.\,.\,.\,.\,.\,.\,.\,.\,.\,.\,.} & $7{,}503\,8375$\\[.3em]
\makebox[22em][l]{$\log(1+C)$\hfill .\,.\,.\,.\,.\,.\,.\,.\,.\,.} & $0{,}000\,0002$\\[.3em]
\makebox[22em][l]{$\mathrm{C.}\log(1-\tfrac{4}{5}T)$\hfill .\,.\,.\,.\,.\,.} & $0{,}001\,1099$\\[.3em]
\hline
\makebox[22em][l]{$\log A$\hfill .\,.\,.\,.\,.\,.\,.\,.\,.\,.\,.\,.\,.\,.} & $7{,}504\,9476$
\end{tabular}
\end{center}
\begin{center}
\begin{tabular}{@{}ll@{}}
\makebox[22em][l]{$\log\dfrac{2B(1+9e)}{15k}\left(\dfrac{q}{e-1}\right)^{\frac{3}{2}}$\hfill .\,.\,.} & $2{,}884\,3582$\\[.4em]
\makebox[22em][l]{$\log A^{\frac{3}{2}}$\hfill .\,.\,.\,.\,.\,.\,.\,.\,.\,.\,.\,.\,.} & $6{,}257\,4214$\\[.3em]
\hline
$\log 0{,}138605$\hfill $=$ & $9{,}141\,7796$
\end{tabular}
\end{center}

\SourceLeaf{064}II. Manentibus $e$ et $q$ ut ante, datur $t=65{,}41236$, quaeruntur $v$ et $r$. Invenimus logarithmos constantium
\begin{center}
\begin{tabular}{@{}ll@{}}
$\log\alpha=$ & $9{,}975\,8345$\\[.3em]
$\log\beta=$ & $9{,}025\,1649$\\[.3em]
$\log\gamma=$ & $9{,}980\,7646$.
\end{tabular}
\end{center}
\noindent Porro prodit $\log\alpha t=1{,}791\,4943$, unde per tabulam \textsc{Barker}i valor approximatus ipsius $w=70^0\,31^{\prime}\,44^{\prime\prime}$, atque hinc $A=0{,}052983$. Huic $A$ in tabula nostra respondet $\log B=0{,}000\,0207$; unde $\log\dfrac{\alpha t}{B}=1{,}791\,4736$, valor correctus ipsius $w=70^0\,31^{\prime}\,36^{\prime\prime}{,}86$. Calculi operationes reliquae ita se habent:
\begin{center}
\begin{tabular}{@{}ll@{}}
\makebox[22em][l]{$2\log\operatorname{tang}\tfrac{1}{2}w$\hfill .\,.\,.\,.\,.\,.} & $9{,}698\,9398$\\[.3em]
\makebox[22em][l]{$\log\beta$\hfill .\,.\,.\,.\,.\,.\,.\,.\,.\,.\,.\,.\,.} & $9{,}025\,1649$\\[.3em]
\hline
\makebox[22em][l]{$\log A$\hfill .\,.\,.\,.\,.\,.\,.\,.\,.\,.\,.\,.\,.} & $8{,}724\,1047$\\[.3em]
$A$\hfill $=$ & $0{,}052\,9791\,1$\\[.3em]
\multicolumn{2}{@{}l@{}}{$\log B$ ut ante}\\[.3em]
$C$\hfill $=$ & $0{,}000\,1261$\\[.3em]
$1+\tfrac{4}{5}A+C$\hfill $=$ & $1{,}042\,5094$\\[.3em]
$1-\tfrac{1}{5}A+C$\hfill $=$ & $0{,}989\,5303$
\end{tabular}
\end{center}
\begin{center}
\begin{tabular}{@{}ll@{}}
\makebox[22em][l]{$\log\operatorname{tang}\tfrac{1}{2}w$\hfill .\,.\,.\,.\,.\,.\,.\,.\,.} & $9{,}849\,4699$\\[.3em]
\makebox[22em][l]{$\log\gamma$\hfill .\,.\,.\,.\,.\,.\,.\,.\,.\,.\,.\,.\,.\,.\,.} & $9{,}980\,7646$\\[.3em]
\makebox[22em][l]{$\tfrac{1}{2}\mathrm{C.}\log(1+\tfrac{4}{5}A+C)$\hfill .\,.} & $9{,}990\,9600$\\[.3em]
\hline
\noalign{\vskip .3em}
\makebox[22em][l]{$\log\operatorname{tang}\tfrac{1}{2}v$\hfill .\,.\,.\,.\,.\,.\,.\,.\,.} & $9{,}821\,1945$\\[.3em]
$\tfrac{1}{2}v$\hfill $=$ & $33^0\,31^{\prime}\,29^{\prime\prime}{,}98$\\[.3em]
$v$\hfill $=$ & $67\quad2\quad59{,}96$\\[.7em]
\makebox[22em][l]{$\log q$\hfill .\,.\,.\,.\,.\,.\,.\,.\,.\,.\,.\,.\,.\,.\,.} & $0{,}020\,1657$\\[.3em]
\makebox[22em][l]{$2\,\mathrm{C.}\log\cos\tfrac{1}{2}v$\hfill .\,.\,.\,.\,.\,.\,.} & $0{,}158\,0378$\\[.3em]
\makebox[22em][l]{$\log(1+\tfrac{4}{5}A+C)$\hfill .\,.\,.\,.\,.} & $0{,}018\,0800$\\[.3em]
\makebox[22em][l]{$\mathrm{C.}\log(1-\tfrac{1}{5}A+C)$\hfill .\,.\,.} & $0{,}004\,5709$\\[.3em]
\hline
\makebox[22em][l]{$\log r$\hfill .\,.\,.\,.\,.\,.\,.\,.\,.\,.\,.\,.\,.\,.} & $0{,}200\,8544$
\end{tabular}
\end{center}
\noindent Quae supra (art. 26) inveneramus $v=67^0\,2^{\prime}\,59^{\prime\prime}{,}78$, $\log r=0{,}200\,8541$, minus exacta sunt, proprieque evadere debuisset $v=67^0\,3^{\prime}\,0^{\prime\prime}{,}00$, quo valore supposito valor ipsius $t$ per tabulas maiores fuerat computatus.
\begin{center}
\rule{7em}{.4pt}
\end{center}
